\section{Experimental Design}
\label{sec:design}

\subsection{Research questions}
\label{sec:rqs}

Six questions organise what follows: three answered by the pre-registered family, three by evidence
labelled exploratory wherever it appears. They expound a design frozen before collection --- no
hypothesis, test or comparison is added, and the family stays at $m=10$ (\S\ref{sec:family}).

\begin{table}[t]
\centering
\caption{The six questions, what answers each, and its status. Confirmatory items map onto
hypotheses frozen before collection; exploratory items are reported as such throughout.}
\label{tab:rqs}
\small
\begin{tabular}{@{}p{0.07\linewidth}p{0.55\linewidth}p{0.30\linewidth}@{}}
\toprule
 & question & answered by \\
\midrule
RQ1 & By how much does the transport of a tool call move the measured score, at fixed model
      and infrastructure? & H1, tests T1--T4 \emph{(confirmatory)} \\
RQ2 & By how much does the serving endpoint move it, at fixed model and transport? &
      H2, tests T5--T8 \emph{(confirmatory)} \\
RQ3 & Is the transport's effect constant across models, and do transport and endpoint
      interact? & H3 and H4, tests T9--T10 \emph{(confirmatory)} \\
RQ4 & By which mechanisms does an endpoint remove a model from an evaluation before it
      produces a score? & the census, \S\ref{sec:census} \emph{(exploratory)} \\
RQ5 & How much of the transport difference is the loss of call batching rather than the
      change of format? & the constrained-native arm, \S\ref{sec:rq5}
      \emph{(exploratory)} \\
RQ6 & Does a power calibration transfer between apparatuses that share model, corpus and
      metric? & \S\ref{sec:results} \emph{(exploratory)} \\
\bottomrule
\end{tabular}
\end{table}

RQ4 through RQ6 were not all asked at the outset: RQ4 was pre-registered as an exploratory census,
while RQ5 and RQ6 are questions the apparatus raised while it ran. What protects them is not that they
were anticipated but that each is answered by a quantity a reader can recompute --- verbatim endpoint
responses, an arm with a stopping criterion fixed in advance, and two calibrations from two
apparatuses.

\subsection{The grid}

Four models, each served by both Databricks and Amazon Bedrock under two transports: sixteen cells. We
name the platforms rather than abstract them --- the venue is single-blind, and the mechanisms of
\S\ref{sec:census} quote the endpoints verbatim, so a reader recognises the platform from the message
whatever we call it and can re-issue each probe against the platform that produced it. Every cell runs
the same 45 held-out binaries eight times, so the design plans 5{,}760 trials while the primary
re-collection delivers 5{,}809 and the replication batch 5{,}805. The three numbers differ for one
reason: a run returning an infrastructure failure is logged as a record and \emph{re-run}, so a cell can
end with more valid measurements than the 360 it planned. The surplus --- 49 above plan and 45 --- is not
a bonus but the count of runs an endpoint refused or dropped, which \S\ref{sec:census} treats as the
subject rather than as attrition, and one cell ends \emph{below} plan at 44 binaries of 45 for the reason
given in \S\ref{sec:ottavo}. Turn budget, temperature and output cap are identical in every cell and
recorded per row rather than assumed.

\begin{table}[h]
\caption{The grid. Every cell is the same corpus under a different apparatus.}
\label{tab:grid}
\centering
\small
\begin{tabular}{ll}
\toprule
models & 4 (two open-weight, two proprietary) \\
infrastructures & 2 (Databricks, Amazon Bedrock) \\
transports & 2 (native, textual) \\
cells & $4\times2\times2 = 16$ \\
binaries per cell & 45, held out, frozen by hash \\
runs per binary & 8 \\
trials & 5{,}760 planned, 5{,}809 measured \\
turn budget & 12 \\
temperature & 0.0 \\
\bottomrule
\end{tabular}
\end{table}

\subsection{Why we say \emph{nominal} model}

Throughout we write \emph{nominal model} rather than \emph{the same weights}. Two managed
clouds advertising the same open-weights model may serve different quantisations, different
runtime versions, or route to a different variant under load, and neither exposes enough to
settle it. Verifying byte-identity is not possible from the outside, so we do not claim it: what
we hold fixed is the model \emph{identity a user selects}, which is exactly the thing a
published evaluation names in its table.

This is not a hedge, it is the object of study. If two endpoints selling the same name behave
differently, that difference belongs in the method section of every paper that names the model
and not the endpoint --- whether the cause is the weights, the runtime, or the routing.

\subsection{Hypotheses, and why two of them are disjunctive}

We ask whether the transport moves the score at fixed infrastructure (H1), whether the
infrastructure moves it at fixed transport (H2), whether the transport's effect is constant
across models (H3), and whether transport and infrastructure interact (H4).

This matters for how the result is read. An interval that excludes effects above some
magnitude is not the same claim as no difference, and we never write the second when the data
support only the first.

\subsection{The test family, and why its size does not shrink}
\label{sec:family}

Ten tests were pre-registered as one family with Holm step-down. The family size is fixed. It
does not shrink when a cell turns out to be inconclusive, uninteresting, or unexecutable,
because removing a member after seeing the data lowers the thresholds of every survivor --- a
correction that grows more permissive exactly as one learns which tests one would like to
pass.

One consequence is deliberate and worth naming: the family includes contrasts nobody
disputes, and those contrasts consume correction budget. That is the price of fixing $m$
before the data exist, and it is cheaper than the alternative.

\subsection{Six defaults inherited from the preceding chapter}

Six constraints are inherited from the preceding chapter, each traceable to a defect that produced a
clean and false number: unwrapped candidate source, functions in native order, symbol tables stripped,
a reasoning-only turn not recorded as a zero, a price guard that refuses \emph{before} the provider
call, and the token-limit parameter sent under the one name the endpoint accepts. None is a discovery
--- each is a default a competent implementer would set, which is what makes the list worth publishing
--- and the artifact states each with the defect it answers.
